% =====================================================================
%  Parte 4 — Reducciones y NP-completitud
%  Fragmento: se incluye con \input desde el documento maestro.
% =====================================================================

% ---------- macros locales de esta parte ----------
\newcommand{\npdato}[1]{{\sans\bfseries\color{acento}#1.}\enspace}
\newcommand{\npsi}{\emph{sí}}
\newcommand{\npno}{\emph{no}}
\newcommand{\npIS}{\prob{Independent Set}}
\newcommand{\npVC}{\prob{Vertex Cover}}
\newcommand{\npSC}{\prob{Set Cover}}
\newcommand{\npSP}{\prob{Set Packing}}
\newcommand{\npCL}{\prob{Clique}}
\newcommand{\npSAT}{\prob{sat}}
\newcommand{\npTSAT}{\prob{3-sat}}
\newcommand{\npSS}{\prob{Subset Sum}}
\newcommand{\npTDM}{\prob{3d-Matching}}
\newcommand{\npTSP}{\prob{tsp}}
\newcommand{\npneg}[1]{\overline{#1}}

\parte{np}{Reducciones y \texorpdfstring{$\NP$}{NP}-completitud}%
{\begin{gathered}X\ \NP\text{-completo},\ \ X\redp Y,\ \ Y\in\NP\\ \Longrightarrow\ \ Y\ \NP\text{-completo}\end{gathered}}%
{Para probar que un problema nuevo es difícil no se lo resuelve: se muestra que, si se lo pudiera resolver rápido, también se resolvería rápido uno que ya sabemos que es difícil. Todo el capítulo es aprender a hacer esa transformación y a escribirla bien.}

% =====================================================================
\section{Conceptos: qué es reducir y qué es \texorpdfstring{$\NP$}{NP}}
% =====================================================================

\subsection{Problemas de decisión y de optimización}

En toda esta parte los problemas son \emph{de decisión}: la respuesta es \npsi\ o \npno. Formalmente, un problema $X$ es el conjunto de tiras (instancias) cuya respuesta es \npsi, y un algoritmo $A$ resuelve $X$ si para toda entrada $s$ vale $A(s)=\text{\npsi}\iff s\in X$. El algoritmo es polinomial si su tiempo es $\Oh(p(|s|))$ para algún polinomio $p$, donde $|s|$ es el largo en bits de la entrada.

Los problemas de optimización se convierten en problemas de decisión agregando un umbral $k$: si se maximiza, la pregunta es <<¿existe una solución de valor \emph{al menos} $k$?>>; si se minimiza, <<¿existe una de valor \emph{a lo sumo} $k$?>>. La versión de decisión no es más difícil que la de optimización (se resuelve la de optimización y se compara con $k$), y tampoco es mucho más fácil: con búsqueda binaria sobre $k$ alcanzan $\Oh(\log n)$ llamadas a la versión de decisión para obtener el valor óptimo. Por eso, para hablar de dificultad, alcanza con mirar la versión de decisión.

\subsection{Reducciones polinomiales}

\begin{definicion}[Kleinberg--Tardos, sección 8.1]
$Y\redp X$ (<<$Y$ se reduce polinomialmente a $X$>>) si cualquier instancia de $Y$ se puede resolver con una cantidad polinomial de pasos de cómputo más una cantidad polinomial de llamadas a una caja negra que resuelve $X$.
\end{definicion}

En la práctica, y en todas las soluciones oficiales, la reducción tiene una forma más restringida: una \emph{transformación} $f$ que toma una instancia $s$ de $Y$, corre en tiempo polinomial en $|s|$, y cumple
\[
  s \text{ es instancia \npsi\ de } Y \iff f(s) \text{ es instancia \npsi\ de } X .
\]
Después se llama una sola vez a la caja negra de $X$ con $f(s)$ y se devuelve su respuesta. Notá que $|f(s)|$ queda polinomial en $|s|$ automáticamente, porque $f$ no tiene tiempo de escribir algo más largo.

\begin{resultado}[Propiedades que se citan en las soluciones]
\begin{enumerate}[label=(\roman*)]
  \item Si $Y\redp X$ y $X$ se resuelve en tiempo polinomial, entonces $Y$ también.
  \item Si $Y\redp X$ y $Y$ no se puede resolver en tiempo polinomial, entonces $X$ tampoco.
  \item Transitividad: si $Z\redp Y$ e $Y\redp X$, entonces $Z\redp X$ (se componen las transformaciones; la composición de polinomios es un polinomio).
  \item Con la definición del libro, si $Y\in\Pclass$ entonces $Y\redp X$ para cualquier $X$: se resuelve $Y$ directamente y se llama cero veces a la caja negra.
\end{enumerate}
\end{resultado}

La lectura que hay que tener automatizada: $Y\redp X$ quiere decir <<$X$ es al menos tan difícil como $Y$>>, o <<si tuviera un algoritmo para $X$, resolvería $Y$>>. Para probar que un problema \emph{nuevo} es difícil, el problema \emph{conocido} va a la izquierda: se transforman instancias del conocido en instancias del nuevo.

\begin{problema}{Ejemplo base: Independent Set $\redp$ Vertex Cover}
En un grafo $G=(V,E)$, $S$ es un conjunto independiente si y solo si $V\setminus S$ es un cubrimiento de vértices: una arista con los dos extremos en $S$ es exactamente una arista que $V\setminus S$ no cubre.

\npdato{Transformación} dada $(G,k)$ de \npIS\ con $n$ vértices, se construye $(G,\,n-k)$ de \npVC. Es $\Oh(1)$ si $n$ está en la representación, $\Oh(n+m)$ si hay que copiar el grafo.

\npdato{Equivalencia} si $S$ es independiente con $|S|\ge k$, $V\setminus S$ es un cubrimiento con $|V\setminus S|\le n-k$; recíprocamente, si $C$ es un cubrimiento con $|C|\le n-k$, $V\setminus C$ es independiente de tamaño $\ge k$. La misma idea da $\npVC\redp\npIS$: los dos problemas son equivalentes.
\end{problema}

Esta transformación se pidió tal cual, sin demostración, como pregunta corta.\visto{2P 2025 Ej. 1(e): IS $\redp$ VC}

\subsection{Certificadores, \texorpdfstring{$\Pclass$}{P} y \texorpdfstring{$\NP$}{NP}}

$\Pclass$ es la clase de los problemas de decisión que tienen un algoritmo polinomial. $\NP$ es la clase de los problemas para los que una respuesta \npsi\ se puede \emph{verificar} rápido si alguien nos da una prueba. La definición que se pide de memoria es la siguiente.

\begin{definicion}[certificador eficiente]
Un algoritmo $B$ es un certificador eficiente para el problema $X$ si
\begin{enumerate}[label=(\arabic*)]
  \item $B$ recibe dos argumentos, una instancia $s$ y un certificado $t$;
  \item el tiempo de ejecución de $B$ es polinomial en $|s|+|t|$;
  \item existe un polinomio $p$ tal que, para toda tira $s$: $s\in X$ si y solo si existe $t$ con $|t|\le p(|s|)$ y $B(s,t)=\text{\npsi}$.
\end{enumerate}
\end{definicion}

\begin{definicion}
$\NP$ es la clase de los problemas de decisión que tienen un certificador eficiente.
\end{definicion}

La condición (3) tiene dos direcciones y en el parcial hay que probar las dos: si $s$ es \npsi\ \emph{existe} un certificado corto que $B$ acepta, y si $B$ acepta \emph{algún} $t$, entonces $s$ es \npsi. Se pidieron las tres definiciones (certificador, $\NP$, $\NP$-completo) como pregunta teórica pura.\visto{Examen feb. 2019 Ej. 3: definiciones}

\begin{itemize}
  \item $\Pclass\subseteq\NP$: si $A$ resuelve $X$ en tiempo polinomial, el certificador $B(s,t)=A(s)$ ignora $t$ (se toma $t$ vacío).
  \item Todo problema de $\NP$ se resuelve en tiempo exponencial: se prueban los $2^{p(|s|)}$ certificados posibles.
  \item $\NP$ \emph{no} quiere decir <<no polinomial>>; quiere decir <<polinomial no determinístico>>, o sea, <<verificable en tiempo polinomial>>. Lo que no se sabe es si $\Pclass=\NP$.
\end{itemize}

El certificado canónico es el de \npSAT: una asignación de valores de verdad a las variables. El certificador evalúa cada cláusula con esa asignación y responde \npsi\ si todas se satisfacen; es lineal en el tamaño de la fórmula. Esta descripción se pidió en dos exámenes recientes, sin demostración de corrección.\visto{Examen dic. 2025 Ej. 1(e): certificado de SAT}

Otros certificados típicos: para \npIS, \npVC\ o \npSC, el conjunto elegido (como arreglo de bits o lista de índices); para ciclo hamiltoniano, la lista de vértices en el orden del ciclo; para 3-coloreo, el color de cada vértice; para \npSS, el subconjunto.\visto{Examen jul. 2026 Ej. 1(f): certificado de SAT}

\subsection{\texorpdfstring{$\NP$}{NP}-completo, \texorpdfstring{$\NP$}{NP}-difícil y Cook--Levin}

\begin{definicion}
Un problema $X$ es $\NP$-completo si (a) $X\in\NP$ y (b) para todo $Y\in\NP$ vale $Y\redp X$. Si se cumple (b) pero no necesariamente (a), $X$ es $\NP$-difícil.
\end{definicion}

\begin{resultado}[Los dos hechos que se usan siempre]
\begin{enumerate}[label=(\roman*)]
  \item Si $X$ es $\NP$-completo, entonces $X$ se resuelve en tiempo polinomial si y solo si $\Pclass=\NP$. (Si $X\in\Pclass$, todo $Y\in\NP$ cumple $Y\redp X$ y por la propiedad (i) $Y\in\Pclass$. Si $\Pclass=\NP$, como $X\in\NP$ entonces $X\in\Pclass$.)
  \item Si $Z$ es $\NP$-completo, $X\in\NP$ y $Z\redp X$, entonces $X$ es $\NP$-completo. (Para $Y\in\NP$: $Y\redp Z$ por ser $Z$ completo, $Z\redp X$ por hipótesis, y por transitividad $Y\redp X$.)
\end{enumerate}
\end{resultado}

El hecho (ii) es toda la estrategia: nunca se prueba (b) a mano para todos los $Y$; se elige un $Z$ conocido y se prueba una sola reducción $Z\redp X$. La definición de $\NP$-completo y la explicación de cómo usarla se pidieron como pregunta corta.\visto{Examen dic. 2025 Ej. 1(f): definir $\NP$-completo}

Hace falta un primer problema $\NP$-completo para arrancar la cadena. Ese es el teorema de Cook--Levin. El libro lo prueba para \prob{Circuit Satisfiability}: dado un circuito booleano con compuertas \textsc{and}, \textsc{or}, \textsc{not} y algunas entradas libres, ¿hay valores para las entradas que hagan que la salida sea 1? La idea de la prueba: si $X\in\NP$ tiene certificador $B$, para cada tamaño de entrada se puede <<desenrollar>> el cómputo polinomial de $B$ en un circuito de tamaño polinomial; se fijan las entradas correspondientes a $s$ y se dejan libres las de $t$. El circuito es satisfacible si y solo si existe $t$ con $B(s,t)=\text{\npsi}$. Después, $\prob{Circuit-sat}\redp\npTSAT$ con una variable por compuerta y unas pocas cláusulas por compuerta que fuerzan su valor; con eso \npTSAT\ queda $\NP$-completo y de ahí sale el resto del árbol.

En el curso, \npSAT, \npTSAT\ y todos los problemas del catálogo de la sección~\ref{sec:np-catalogo} se dan por conocidos $\NP$-completos: se pueden usar sin demostrarlo, citando el libro.

\subsection{Más allá de \texorpdfstring{$\NP$}{NP}: co-NP y PSPACE}

En los parciales relevados no se preguntó nada de esto, pero conviene tener la idea para no confundirse en un verdadero o falso.

\begin{itemize}
  \item $X\in\coNP$ si su complemento $\overline{X}$ (las instancias \npno) está en $\NP$. Ejemplos: <<¿la fórmula es insatisfacible?>>, <<¿el grafo \emph{no} tiene ciclo hamiltoniano?>>. No se sabe si $\NP=\coNP$. Como $\Pclass$ es cerrada por complemento, si $\NP\ne\coNP$ entonces $\Pclass\ne\NP$.
  \item $\Pclass\subseteq\NP\cap\coNP$. Un problema en $\NP\cap\coNP$ tiene una <<buena caracterización>>: hay certificado corto tanto para el \npsi\ como para el \npno. Ejemplo del curso: <<¿hay un flujo de valor al menos $v$?>>; un flujo certifica el \npsi\ y un corte de capacidad menor a $v$ certifica el \npno.
  \item $\mathsf{PSPACE}$ es la clase de problemas que se resuelven con memoria polinomial. Vale $\Pclass\subseteq\NP\subseteq\mathsf{PSPACE}$ (se prueban todos los certificados reutilizando el mismo espacio). El problema completo típico es \prob{qsat} (fórmulas con cuantificadores alternados), que modela juegos de dos jugadores y planificación.
\end{itemize}

\subsection{Preguntas conceptuales de verdadero o falso}

Las preguntas de este tipo se contestan con una tabla mental de implicaciones. Lo único que importa es de qué lado de $\redp$ está cada problema.

\begin{resultado}[Qué se puede deducir de $Y\redp X$]
\begin{itemize}
  \item $X\in\Pclass\Rightarrow Y\in\Pclass$. \quad $Y\notin\Pclass\Rightarrow X\notin\Pclass$.
  \item $Y\in\Pclass$: no dice nada de $X$. \quad $X$ $\NP$-completo: no dice nada de $Y$ (salvo si $Y\in\NP$, que ya se sabía).
  \item $Y$ $\NP$-completo y $X\in\NP$: $X$ es $\NP$-completo. \quad $Y$ $\NP$-completo y $X\in\Pclass$: $\Pclass=\NP$.
  \item Para todo $Y\in\NP$ y todo $Z$ $\NP$-completo: $Y\redp Z$, siempre. En particular todo problema de $\Pclass$ se reduce a cualquier $\NP$-completo.
  \item Si $\Pclass\ne\NP$: ningún $\NP$-completo está en $\Pclass$, y ningún $\NP$-completo se reduce a un problema de $\Pclass$.
\end{itemize}
\end{resultado}

\begin{problema}{Sí, no o no se sabe (2P 2017, Ej. 4)}
Sea \prob{Interval Scheduling} en versión de decisión: dada una colección de intervalos y una cota $k$, ¿hay al menos $k$ intervalos que no se superponen? Para cada afirmación, responder <<sí>>, <<no>> o <<no se sabe; si se supiera, se resolvería $\Pclass=\NP$>>.

\npdato{(a) ¿Interval Scheduling $\redp$ Vertex Cover?} Sí. Interval Scheduling está en $\Pclass$ (greedy por tiempo de finalización), así que se resuelve en tiempo polinomial con una caja negra de \npVC\ llamándola cero veces. Otra justificación: Interval Scheduling está en $\NP$ y todo problema de $\NP$ se reduce a \npVC, que es $\NP$-completo.

\npdato{(b) ¿Independent Set $\redp$ Interval Scheduling?} No se sabe; es equivalente a $\Pclass=\NP$. Si fuera cierto, como Interval Scheduling está en $\Pclass$, \npIS\ estaría en $\Pclass$ y por ser $\NP$-completo $\Pclass=\NP$. Recíprocamente, si $\Pclass=\NP$, \npIS\ se resuelve en tiempo polinomial y se reduce a cualquier cosa.
\end{problema}

Esta pregunta (la misma, con las mismas respuestas) aparece en la práctica de segundo parcial del curso.\visto{2P 2017 Ej. 4: sí / no / no se sabe}

\begin{problema}{Cuatro afirmaciones (examen jul. 2018, Pregunta 2)}
Sean $X\in\Pclass$, $Y\in\NP$ y $Z$ $\NP$-completo. Decir si cada afirmación es siempre verdadera, siempre falsa o no se puede afirmar; en el último caso, qué tendría que cumplirse para que sea verdadera.

\npdato{1. $Y\redp Z$} Siempre verdadera: es la definición de $\NP$-completo.

\npdato{2. $X\redp Z$} Siempre verdadera: $X\in\Pclass\subseteq\NP$ y se aplica el punto 1.

\npdato{3. $Z\redp Y$} No se puede afirmar. Es verdadera exactamente cuando $Y$ es $\NP$-completo: si vale, $Y\in\NP$ y es al menos tan difícil como un $\NP$-completo.

\npdato{4. $Z\redp X$} No se puede afirmar. Si fuera verdadera, $Z$ se resolvería en tiempo polinomial y entonces $\Pclass=\NP$; recíprocamente, si $\Pclass=\NP$ vale (propiedad (iv)). Es verdadera si y solo si $\Pclass=\NP$.
\end{problema}

Ninguna de las dos preguntas pide demostraciones: piden la justificación de una o dos oraciones basada en las propiedades de arriba.\visto{Examen jul. 2018 Preg. 2: X, Y, Z}

Para practicar, otras afirmaciones del mismo estilo (no salieron tal cual, pero se contestan con la misma tabla):

\begin{enumerate}[label=\alph*)]
  \item <<Si $Y\redp X$ e $Y\in\Pclass$, entonces $X\in\Pclass$>>. Falso en general: $Y$ fácil no dice nada de $X$ (todo problema de $\Pclass$ se reduce a \npSAT).
  \item <<Si $X\redp\npTSAT$, entonces $X$ es $\NP$-completo>>. Falso: prueba que $X$ es \emph{a lo sumo} tan difícil como \npTSAT. La dirección correcta es $\npTSAT\redp X$ (y además $X\in\NP$).
  \item <<\npSS\ tiene un algoritmo $\Oh(nW)$, por lo tanto $\Pclass=\NP$>>. Falso: $\Oh(nW)$ es exponencial en la cantidad de bits de $W$.
  \item <<Si $\Pclass\ne\NP$, todo problema de $\NP$ que no está en $\Pclass$ es $\NP$-completo>>. Falso: si $\Pclass\ne\NP$ existen problemas intermedios (teorema de Ladner). No hace falta saberlo para el parcial, pero sí no afirmar lo contrario.
  \item <<Dos problemas $\NP$-completos se reducen uno al otro>>. Verdadero: ambos están en $\NP$ y cada uno es completo.
  \item <<Si un problema $\NP$-completo está en $\coNP$, entonces $\NP=\coNP$>>. Verdadero (sección 8.9 del libro), aunque no se usa en los parciales.
\end{enumerate}

% =====================================================================
\section{Catálogo de problemas \texorpdfstring{$\NP$}{NP}-completos del libro}
\label{sec:np-catalogo}
% =====================================================================

Estos son los problemas que se pueden usar como <<problema de referencia>> sin demostrar que son $\NP$-completos. Para cada uno: la versión de decisión, el certificado y la reducción con la que el libro lo prueba. En los parciales nunca se pidió reproducir una reducción del libro con gadgets complicados, pero sí reconocer los problemas y saber de cuál viene cada uno.

\subsection{Empaquetamiento y cobertura}

\npdato{\npIS} dado un grafo $G=(V,E)$ y un natural $k$, ¿existe $S\subseteq V$ con $|S|\ge k$ y sin aristas entre vértices de $S$? Certificado: $S$. Se prueba desde \npTSAT\ (ver abajo).

\npdato{\npVC} dado $G$ y $k$, ¿existe $C\subseteq V$ con $|C|\le k$ tal que toda arista tiene al menos un extremo en $C$? Certificado: $C$. Se prueba desde \npIS\ con $(G,k)\mapsto(G,n-k)$.

\npdato{\npSC} dado un conjunto $U$ de $n$ elementos, subconjuntos $S_1,\dots,S_m\subseteq U$ y $k$, ¿hay a lo sumo $k$ de los $S_i$ cuya unión es $U$? Certificado: los índices elegidos. Se prueba desde \npVC: $U=E$ y, para cada vértice $v$, $S_v=\{\text{aristas incidentes a }v\}$, con el mismo $k$. Un conjunto de vértices cubre todas las aristas si y solo si los $S_v$ correspondientes cubren $U$.

\npdato{\npSP} dado $U$, subconjuntos $S_1,\dots,S_m$ y $k$, ¿hay al menos $k$ de los $S_i$ disjuntos dos a dos? Se prueba desde \npIS\ con la misma construcción: $U=E$, $S_v$ las aristas incidentes a $v$. Dos $S_v$ son disjuntos si y solo si $u$ y $v$ no son adyacentes.

\npdato{\npCL} dado $G$ y $k$, ¿hay $k$ vértices adyacentes dos a dos? Se prueba desde \npIS\ con el grafo complemento $\overline{G}$ (misma $V$, arista donde $G$ no tiene) y el mismo $k$: $S$ es independiente en $G$ si y solo si es clique en $\overline{G}$. Construir $\overline G$ es $\Oh(n^2)$.

\npdato{\prob{Hitting Set}} dado $A$, subconjuntos $B_1,\dots,B_m$ de $A$ y $k$, ¿hay $H\subseteq A$ con $|H|\le k$ que toca a todos los $B_i$? Es el <<transpuesto>> de \npSC\ (se eligen elementos en lugar de conjuntos). Desde \npVC: $A=V$ y un $B_i=\{u,v\}$ por cada arista.

\subsection{Satisfacibilidad}

\npdato{\npSAT} dadas variables $x_1,\dots,x_n$ y cláusulas $C_1,\dots,C_k$ (cada una disyunción de términos $x_i$ o $\npneg{x_i}$), ¿hay una asignación que satisface todas las cláusulas? \npTSAT\ es el caso en que cada cláusula tiene exactamente tres términos sobre variables distintas. Certificado: la asignación. \npTSAT\ se prueba desde \prob{Circuit-sat} (Cook--Levin).

Una reducción que viene bien conocer porque es la base de los ejercicios de <<$k$-\prob{sat}>>: $\npSAT\redp\npTSAT$. Una cláusula larga $(t_1\vee t_2\vee\dots\vee t_r)$ se parte con variables nuevas $y_j$:
\[
(t_1\vee t_2\vee y_1)\wedge(\npneg{y_1}\vee t_3\vee y_2)\wedge\dots\wedge(\npneg{y_{r-3}}\vee t_{r-1}\vee t_r).
\]
Si algún $t_i$ es verdadero, se eligen los $y_j$ para satisfacer el resto de la cadena; si todos los $t_i$ son falsos, la cadena obliga $y_1=1$, luego $y_2=1$, \dots, y la última cláusula falla. Las cláusulas cortas se completan con variables nuevas forzadas.

\begin{problema}{3-SAT $\redp$ Independent Set (libro, sección 8.2)}
\npdato{Construcción} dada una instancia de \npTSAT\ con $k$ cláusulas, se arma un grafo con $3k$ vértices: un triángulo por cláusula, con un vértice por cada término. Además se une con una arista cada par de vértices que representan términos opuestos ($x_i$ en una cláusula y $\npneg{x_i}$ en otra). Se pregunta por un independiente de tamaño $k$.

\npdato{Ida} si hay asignación que satisface, se elige en cada triángulo un término verdadero. Son $k$ vértices, no hay dos del mismo triángulo y no hay dos opuestos (no pueden ser ambos verdaderos).

\npdato{Vuelta} un independiente de tamaño $k$ tiene exactamente un vértice por triángulo y no contiene opuestos. Se pone verdadero cada término elegido (no hay conflicto) y el resto de las variables arbitrarias: cada cláusula tiene un término verdadero.
\end{problema}

\subsection{Secuenciación}

\npdato{Ciclo hamiltoniano} dado un grafo (dirigido o no), ¿hay un ciclo simple que pasa por todos los vértices? Certificado: los vértices en orden. El libro lo prueba para grafos dirigidos desde \npTSAT; la versión no dirigida sale de la dirigida partiendo cada vértice $v$ en un camino $v_{\text{in}}-v_{\text{mid}}-v_{\text{out}}$.

\npdato{Camino hamiltoniano} igual, pero un camino simple (no hace falta volver). Desde ciclo hamiltoniano dirigido: se elige un vértice $v$ y se lo parte en $v'$, que se queda con las aristas salientes de $v$, y $v''$, que se queda con las entrantes. Como $v'$ no tiene entrantes y $v''$ no tiene salientes, todo camino hamiltoniano va de $v'$ a $v''$, y eso es lo mismo que un ciclo hamiltoniano por $v$ en el grafo original.

\npdato{\npTSP} dadas $n$ ciudades, distancias $d(u,v)$ y una cota $D$, ¿hay un recorrido que visita cada ciudad una vez, vuelve al inicio y tiene largo total $\le D$? Desde ciclo hamiltoniano: $d(u,v)=1$ si $(u,v)$ es arista y $2$ si no, con $D=n$. Un recorrido de largo $n$ usa solo aristas del grafo.

\begin{problema}{3-SAT $\redp$ ciclo hamiltoniano dirigido (libro, sección 8.5)}
\npdato{Gadget de variable} por cada variable $x_i$ una fila de nodos $v_{i,1},\dots,v_{i,b}$ ($b=3k+3$) con aristas en ambos sentidos entre vecinos. Recorrer la fila de izquierda a derecha codifica $x_i=1$; de derecha a izquierda, $x_i=0$. Los extremos de cada fila se conectan con los dos extremos de la fila siguiente, y se agregan un nodo inicial $s$, uno final $t$ y la arista $(t,s)$. Sin cláusulas, hay exactamente $2^n$ ciclos hamiltonianos, uno por asignación.

\npdato{Gadget de cláusula} por cada cláusula $C_j$ un nodo $c_j$. Si $x_i$ aparece en $C_j$ se agregan $v_{i,3j}\to c_j\to v_{i,3j+1}$ (un desvío que solo se puede tomar yendo de izquierda a derecha); si aparece $\npneg{x_i}$, se agregan $v_{i,3j+1}\to c_j\to v_{i,3j}$ (solo yendo de derecha a izquierda).

\npdato{Equivalencia} una asignación que satisface da un ciclo: se recorren las filas según la asignación y cada $c_j$ se visita por el desvío de un término verdadero. Recíprocamente, en un ciclo hamiltoniano cada $c_j$ se visita por un desvío que vuelve a la misma fila, y el sentido de esa fila hace verdadero al término correspondiente.
\end{problema}

\subsection{Partición: 3D-Matching}

\npdato{\npTDM} dados conjuntos disjuntos $X,Y,Z$ de tamaño $n$ cada uno y un conjunto $T\subseteq X\times Y\times Z$ de ternas, ¿hay $n$ ternas de $T$ que cubren cada elemento exactamente una vez? Es la generalización a tres dimensiones del matching bipartito perfecto (que sí está en $\Pclass$, por flujo). Certificado: las $n$ ternas.

\begin{problema}{3-SAT $\redp$ 3D-Matching (libro, sección 8.6), en una línea por gadget}
\npdato{Variable} un <<anillo>> de $2k$ elementos núcleo y $2k$ puntas, con ternas (núcleo$_j$, núcleo$_{j+1}$, punta$_j$). Para cubrir el núcleo hay que tomar todas las ternas pares o todas las impares: eso deja libres las puntas impares ($x_i=1$) o las pares ($x_i=0$).

\npdato{Cláusula} dos elementos núcleo nuevos y una terna con cada punta libre que corresponde a un término de la cláusula: solo se puede cubrir si algún término es verdadero.

\npdato{Limpieza} gadgets extra, cada uno con ternas con todas las puntas, para cubrir las puntas que sobran. Se reparten los elementos en $X$, $Y$, $Z$ de modo que cada terna tenga uno de cada uno.
\end{problema}

\subsection{Coloreo}

\npdato{$k$-coloreo} dado $G$, ¿se pueden asignar $k$ colores a los vértices de modo que vecinos tengan colores distintos? Para $k=2$ está en $\Pclass$ (es ver si el grafo es bipartito, con BFS). Para $k=3$ ya es $\NP$-completo. Certificado: el color de cada vértice. De $3$-coloreo a $k$-coloreo, $k>3$: se agregan $k-3$ vértices nuevos adyacentes a todos (y entre sí); se llevan $k-3$ colores propios y quedan 3 para el grafo original.

\begin{problema}{3-SAT $\redp$ 3-coloreo (libro, sección 8.7)}
\npdato{Paleta} un triángulo con nodos especiales $T$, $F$ y $B$ (base); sus colores representan verdadero, falso y <<otro>>.

\npdato{Variables} por cada $x_i$, dos nodos $v_i$ y $\npneg{v_i}$ unidos entre sí y con $B$. Así cada uno toma color de $T$ o de $F$, y son opuestos.

\npdato{Cláusulas} por cada cláusula, un subgrafo de seis nodos conectado a sus tres términos y a $T$ y $F$, armado para que se pueda colorear si y solo si alguno de los tres términos tiene el color de $T$. Si los tres son $F$, el nodo de salida del gadget queda forzado a $F$ y choca con una arista a $F$.
\end{problema}

\subsection{Problemas numéricos}

\npdato{\npSS} dados naturales $w_1,\dots,w_n$ y una meta $W$, ¿hay un subconjunto que suma exactamente $W$? Certificado: el subconjunto. Tiene un algoritmo $\Oh(nW)$ que no es polinomial (ver arriba).

\npdato{Mochila (decisión)} dados pesos $w_i$, valores $v_i$, capacidad $W$ y meta $V$, ¿hay un subconjunto de peso $\le W$ y valor $\ge V$? Desde \npSS: $v_i=w_i$ y $V=W$; peso $\le W$ y valor $\ge W$ obligan a sumar exactamente $W$.

\npdato{Secuenciación con release times y deadlines} $n$ trabajos en una sola máquina, el trabajo $i$ dura $t_i$, puede empezar desde $r_i$ y debe terminar antes de $d_i$; ¿hay una planificación sin solapamientos que cumple todo? Desde \npSS: un trabajo por cada $w_i$, con $r_i=0$, $t_i=w_i$, $d_i=S+1$ (donde $S=\sum w_i$), y un trabajo extra con $r=W$, $t=1$, $d=W+1$. El extra ocupa exactamente $[W,W+1]$, así que antes de él tiene que haber trabajos que llenan exactamente $W$ unidades.

\begin{problema}{3D-Matching $\redp$ Subset Sum (libro, sección 8.8)}
\npdato{Construcción} con $|X|=|Y|=|Z|=n$ y $m$ ternas, a cada terna $(x_i,y_j,z_k)$ se le asocia el número que en base $b=m+1$ tiene un 1 en las posiciones $i$, $n+j$ y $2n+k$ y 0 en las demás. La meta es $W=\sum_{r=0}^{3n-1}b^{r}$ (todos los dígitos en 1).

\npdato{Por qué funciona} como hay a lo sumo $m$ números y cada dígito vale 0 o 1, nunca hay acarreo en base $m+1$. Sumar $W$ es tener exactamente un 1 en cada posición, es decir, cubrir cada elemento exactamente una vez. Cada número tiene $3n$ dígitos en base $m+1$, o sea $\Oh(n\log m)$ bits: la transformación es polinomial.
\end{problema}

\subsection{El árbol de reducciones}

Cada flecha $A\to B$ significa $A\redp B$ en la prueba del libro: se transforman instancias de $A$ en instancias de $B$, y $B$ hereda la dificultad.

\begin{center}
\begin{tikzpicture}[
    nodo/.style={draw=acento,fill=acento!7,rounded corners=2pt,font=\sans\footnotesize,
                 align=center,inner sep=3pt,minimum height=0.62cm},
    raiz/.style={nodo,fill=acento!22},
    fl/.style={-{Stealth[length=2mm]},draw=gris,thick},
    x=1cm,y=1cm]
  \node[raiz] (csat) at (6.2,0)    {Circuit-SAT\\[-1pt]{\scriptsize Cook--Levin}};
  \node[raiz] (sat)  at (6.2,-1.5) {SAT $\equiv$ 3-SAT};
  \node[nodo] (is)   at (2.4,-3.1) {Independent Set};
  \node[nodo] (hc)   at (6.1,-3.1) {ciclo\\[-1pt]hamiltoniano};
  \node[nodo] (tdm)  at (8.9,-3.1) {3D-Matching};
  \node[nodo] (col)  at (11.6,-3.1){3-coloreo};
  \node[nodo] (vc)   at (0.75,-4.7){Vertex\\[-1pt]Cover};
  \node[nodo] (sp)   at (2.4,-4.7) {Set\\[-1pt]Packing};
  \node[nodo] (cl)   at (3.95,-4.7){Clique};
  \node[nodo] (hp)   at (5.45,-4.7){camino\\[-1pt]ham.};
  \node[nodo] (tsp)  at (6.9,-4.7) {TSP};
  \node[nodo] (ss)   at (8.9,-4.7) {Subset Sum};
  \node[nodo] (kcol) at (11.6,-4.7){$k$-coloreo};
  \node[nodo] (sc)   at (0.75,-6.3){Set\\[-1pt]Cover};
  \node[nodo] (hs)   at (2.4,-6.3) {Hitting\\[-1pt]Set};
  \node[nodo] (sch)  at (7.95,-6.3){secuenciación con\\[-1pt]release times};
  \node[nodo] (mo)   at (10.25,-6.3){mochila};
  \draw[fl] (csat) -- (sat);
  \draw[fl] (sat) -- (is);  \draw[fl] (sat) -- (hc);
  \draw[fl] (sat) -- (tdm); \draw[fl] (sat) -- (col);
  \draw[fl] (is) -- (vc); \draw[fl] (is) -- (sp); \draw[fl] (is) -- (cl);
  \draw[fl] (vc) -- (sc); \draw[fl] (vc) -- (hs);
  \draw[fl] (hc) -- (hp); \draw[fl] (hc) -- (tsp);
  \draw[fl] (tdm) -- (ss);
  \draw[fl] (ss) -- (sch); \draw[fl] (ss) -- (mo);
  \draw[fl] (col) -- (kcol);
\end{tikzpicture}
\end{center}

\npVC\ e \npIS\ se reducen en ambos sentidos; también \npSAT\ y \npTSAT. \prob{Hitting Set} y \npCL\ no están entre las reducciones que el libro desarrolla en el texto (son ejercicios estándar), pero son tan directas que conviene tenerlas en el árbol: de ahí cuelgan varios ejercicios de examen.

% =====================================================================
\section{Cómo se escribe una demostración de \texorpdfstring{$\NP$}{NP}-completitud}
% =====================================================================

Todas las soluciones oficiales tienen la misma estructura, y desde 2024 la pauta de corrección la explicita punto por punto. Si se sigue la plantilla, el puntaje está casi asegurado aunque la reducción sea la obvia.

\begin{receta}[demostrar que $X$ es $\NP$-completo]
\textbf{0. Plantear.} <<Para probar que $X$ es $\NP$-completo hay que probar (I) $X\in\NP$ y (II) $Y\redp X$ para todo $Y\in\NP$. Para (II) alcanza con probar $Z\redp X$ para un $Z$ $\NP$-completo conocido y usar la transitividad de $\redp$.>>

\textbf{1. $X\in\NP$.}
\begin{itemize}
  \item Describir la instancia $s$ y su representación (matriz de adyacencia, lista de cláusulas, arreglo\dots).
  \item Definir el certificado $t$ (casi siempre: la solución candidata, como arreglo de bits o lista de índices) y decir por qué $|t|\le p(|s|)$.
  \item Escribir el certificador $B(s,t)$ en pseudocódigo: una serie de chequeos que responden \npno\ si alguno falla y \npsi\ al final.
  \item Tiempo de $B$: polinomial en $|s|+|t|$ (contar los ciclos).
  \item Si $s$ es \npsi, existe $t$ (la solución misma) con $B(s,t)=\text{\npsi}$.
  \item Si $B(s,t)=\text{\npsi}$ para algún $t$, cada chequeo garantiza una condición del problema, así que $s$ es \npsi.
\end{itemize}

\textbf{2. $Z\redp X$.}
\begin{itemize}
  \item Tomar una instancia \emph{arbitraria} $s_Z$ del problema conocido $Z$ y describir la instancia $s_X=f(s_Z)$, componente por componente (<<para cada vértice $v_i$ se define\dots>>).
  \item Justificar que $f$ corre en tiempo polinomial en $|s_Z|$ (y por lo tanto $|s_X|$ es polinomial).
  \item Ida: si $s_Z$ es \npsi, construir explícitamente una solución de $s_X$ a partir de la de $s_Z$.
  \item Vuelta: si $s_X$ es \npsi, construir una solución de $s_Z$ a partir de la de $s_X$. (Equivalente: si $s_Z$ es \npno, $s_X$ es \npno.)
  \item Cerrar: <<se resuelve $Z$ con la transformación, que es polinomial, y una sola llamada a un algoritmo que resuelve $X$; por lo tanto $Z\redp X$>>.
\end{itemize}

\textbf{3. Concluir.} <<Como $Z$ es $\NP$-completo, todo $Y\in\NP$ cumple $Y\redp Z$; por transitividad $Y\redp X$. Junto con $X\in\NP$, $X$ es $\NP$-completo.>>
\end{receta}

Cuando la letra da una hipótesis, se usa y no se demuestra de nuevo: <<asuma que existe un certificador eficiente>> resuelve el paso 1 en una oración; <<asuma que si la original es \npsi\ la transformada es \npsi>> elimina la ida. La pauta de 2025 pide explícitamente \emph{no} describir el certificador en ese caso, sino decir que se asume.\visto{2P 2025 Ej. 2: usar las hipótesis dadas}

\begin{trampa}
\begin{itemize}
  \item \textbf{Reducir al revés.} Probar $X\redp Z$ (el nuevo al conocido) no prueba nada sobre la dificultad de $X$. La pauta de 2024 lo califica de <<error grave>>. Chequeo: la instancia que se toma \emph{arbitraria} es la del problema conocido.
  \item \textbf{Olvidar $X\in\NP$.} Sin eso se probó $\NP$-difícil, no $\NP$-completo.
  \item \textbf{Probar una sola dirección.} La ida suele ser obvia; la vuelta es donde está el ejercicio (\prob{Cometa}, \prob{nti4sat}, \prob{casi-sat}).
  \item \textbf{Vuelta sobre instancias arbitrarias de $X$.} La vuelta se prueba solo para instancias de la forma $f(s_Z)$; se puede (y se debe) usar cómo fueron construidas.
  \item \textbf{Transformación que usa la respuesta.} $f$ no puede depender de conocer la solución de $s_Z$ (por ejemplo, <<se toma $k$ igual al tamaño del máximo independiente>>), porque calcularla no es polinomial.
  \item \textbf{Construcción no polinomial.} Escribir números en unario, enumerar todos los subconjuntos o crear un objeto por cada asignación de verdad.
  \item \textbf{Instancias inválidas.} La instancia construida tiene que respetar la definición de $X$: en \prob{4-sat} los cuatro términos son sobre variables distintas, en \prob{Red-Social} la amistad es simétrica, en \prob{Cometa} $d\ge 3$.
  \item \textbf{Certificador que busca.} El certificador verifica la solución que recibe; si para responder tiene que buscar entre muchas opciones, no es polinomial.
  \item \textbf{Certificado sin tamaño.} Hay que decir por qué $|t|$ es polinomial en $|s|$, no solo describirlo.
  \item \textbf{Olvidar el tiempo de la transformación.} Aunque sea una copia, hay que decir que es lineal o polinomial.
\end{itemize}
\end{trampa}

\subsection{Cómo elegir el problema de origen}

La sugerencia de la letra suele decir de dónde reducir. Cuando no la hay, se busca el problema del catálogo que se parece más <<estructuralmente>> al nuevo.

\begin{center}\small
\begin{tabularx}{\textwidth}{@{}>{\raggedright\arraybackslash}X >{\raggedright\arraybackslash}p{3.6cm}@{}}
\toprule
Si el problema nuevo\dots & conviene reducir desde \\
\midrule
elige elementos de un grafo (o de una matriz de <<conflictos>>) que no pueden convivir, y pide al menos $k$ & \npIS \\
pide al menos $k$ elementos todos compatibles entre sí, o un subgrafo completo & \npCL\ (o \npIS\ con complemento) \\
pide a lo sumo $k$ objetos que <<tocan>> todas las restricciones de a pares & \npVC \\
pide a lo sumo $k$ grupos cuya unión cubre todo (o $k$ elementos que tocan todos los grupos) & \npSC\ o \npVC \\
tiene variables booleanas, cláusulas, apuestas o tableros con dos estados por fila & \npSAT\ o \npTSAT \\
pide recorrer todo una vez, un camino o ciclo largo, un orden & camino o ciclo hamiltoniano \\
pide sumar exactamente un número, repartir en dos mitades, cargar una capacidad & \npSS \\
pide asignar etiquetas a vértices con restricciones de vecinos & 3-coloreo \\
pide particionar en ternas & \npTDM \\
\bottomrule
\end{tabularx}
\end{center}

Una observación que ahorra tiempo: casi todos los ejercicios de examen son \emph{disfraces} del problema de referencia, y la transformación es un renombre (vértices $\to$ ubicaciones, aristas $\to$ tareas, cláusulas $\to$ columnas). Si la transformación que se está armando es complicada, probablemente se eligió mal el origen.

\subsection{Los problemas de examen, de un vistazo}

La figura muestra de qué problema conocido se redujo cada ejercicio en la solución oficial.

\begin{center}
\begin{tikzpicture}[
    src/.style={draw=acento,fill=acento!14,rounded corners=2pt,font=\sans\footnotesize\bfseries,
                align=center,inner sep=3pt,minimum height=0.6cm,minimum width=2.6cm},
    hoja/.style={draw=grisclaro,fill=white,font=\sans\scriptsize,align=left,
                 text width=2.65cm,inner sep=3pt,anchor=north},
    nada/.style={hoja,draw=gris,dashed,text=gris},
    fl/.style={-{Stealth[length=2mm]},draw=gris,thick},
    x=1cm,y=1cm]
  \node[src] (is) at (1.6,0)  {Independent Set};
  \node[src] (cl) at (4.8,0)  {Clique};
  \node[src] (vc) at (8.0,0)  {Vertex Cover};
  \node[src] (sc) at (11.2,0) {Set Cover};
  \node[hoja] (his) at (1.6,-0.9) {UOPI \hfill 2P 2018\\ Red-Social \hfill feb.\,2021\\ Cocinero \hfill jul.\,2022\\ RF \hfill dic.\,2022};
  \node[hoja] (hcl) at (4.8,-0.9) {Cometa \hfill dic.\,2023\\ Subgrafo isom. \hfill jul.\,2024};
  \node[hoja] (hvc) at (8.0,-0.9) {PCG \hfill feb.\,2020\\ Test \hfill jul.\,2021\\ Contratar \hfill 2P 2022};
  \node[hoja] (hsc) at (11.2,-0.9){Acceso \hfill 2P 2023\\ UV \hfill 2P 2025};
  \node[src] (sat) at (1.6,-3.6)  {SAT};
  \node[src] (tsat) at (4.8,-3.6) {3-SAT};
  \node[src] (hp) at (8.0,-3.6)   {camino ham.};
  \node[src,draw=gris,fill=grisclaro!60] (otros) at (11.2,-3.6) {sin ejercicios};
  \node[hoja] (hsat) at (1.6,-4.5) {casiSAT \hfill 2P 2020\\ Falsable \hfill feb.\,2022\\ Al menos $k$ \hfill jul.\,2025};
  \node[hoja] (htsat) at (4.8,-4.5){Solitaire \hfill 2P 2019\\ 4-SAT \hfill 2P 2024\\ NTI4Sat \hfill jul.\,2026};
  \node[hoja] (hhp) at (8.0,-4.5)  {camino simple\\ de largo $\ge k$ \hfill 2P 2021};
  \node[nada] (hot) at (11.2,-4.5) {Subset Sum, 3-coloreo, 3D-Matching, TSP};
  \foreach \a/\b in {is/his,cl/hcl,vc/hvc,sc/hsc,sat/hsat,tsat/htsat,hp/hhp}
    \draw[fl] (\a) -- (\b);
\end{tikzpicture}
\end{center}

% =====================================================================
\section{Problemas de examen, agrupados por problema de origen}
% =====================================================================

Para cada ejercicio: enunciado resumido, cómo se prueba que está en $\NP$, la transformación y los argumentos de ida y vuelta. Los certificadores de los ejercicios son todos del mismo tipo (recorrer la solución candidata y chequear las restricciones); se detallan en los primeros y después se resumen.

\subsection{Desde Independent Set}

La señal es <<elegir al menos $k$ cosas sin conflictos de a pares>>. La transformación casi siempre deja el grafo igual (o lo codifica en una matriz) y deja $k$ igual.

El primero es el más directo de todos: \npIS\ con pesos. La solución oficial es un buen modelo de redacción del certificador, con las dos direcciones explícitas.\visto{2P 2018 Ej. 3: UOPI, desde IS}

\begin{problema}{Ubicación óptima de plantas industriales (UOPI)}
\npdato{Problema} grafo $G$ cuyos vértices son ubicaciones con valor entero $w_i$; una arista indica que dos ubicaciones son excluyentes. Dado $t$, ¿hay ubicaciones no excluyentes entre sí cuyo valor total es al menos $t$?

\npdato{En $\NP$} certificado: tira de $n$ bits $c$ (ubicaciones elegidas). $B$ suma los $w_i$ con $c_i=1$ y rechaza si la suma es menor a $t$; después, para cada $i$ con $c_i=1$, recorre la lista de adyacencia de $i$ y rechaza si encuentra un $j$ con $c_j=1$. Es $\Oh(n+m)$. Si $B$ acepta, la selección suma al menos $t$ y no tiene pares excluyentes; si la instancia es \npsi, la tira de esa selección hace aceptar a $B$.

\npdato{Reducción $\npIS\redp\prob{uopi}$} $(G,k)\mapsto(G,\,w_i=1\ \forall i,\,t=k)$, en $\Oh(n)$. Con pesos 1, la suma de un conjunto es su tamaño; un conjunto es independiente si y solo si no tiene exclusiones. Así, hay independiente de tamaño $\ge k$ si y solo si hay selección válida de valor $\ge t$.
\end{problema}

El segundo es \npCL\ con otro nombre, y la reducción desde \npIS\ es el grafo complemento: <<amigos>> son los no adyacentes.\visto{Examen feb. 2021 Ej. 3: Red-Social, desde IS}

\begin{problema}{Red-Social}
\npdato{Problema} personas $P=\{p_1,\dots,p_n\}$ con conjuntos de amigos $A_i$ (relación simétrica) y un natural $k$: ¿hay al menos $k$ personas amigas todas entre sí?

\npdato{Reducción $\npIS\redp\prob{Red-Social}$} dada $(G,k)$: $P=V$; $A_i=\{v_j : j\ne i,\ v_j \text{ no adyacente a } v_i\}$; $k'=k$. La relación es simétrica porque la adyacencia lo es. Tiempo $\Oh(n^3)$ con una búsqueda ingenua (o $\Oh(n^2)$ con matriz).

\npdato{Ida y vuelta} $S$ independiente en $G$ $\iff$ cada par de $S$ es no adyacente $\iff$ cada par de $S$ es amigo en $R$. El tamaño se conserva porque $k'=k$.

\npdato{En $\NP$} certificado: arreglo booleano de $n$ posiciones. $B$ rechaza si hay menos de $k$ elegidos y, para cada par $(p,q)$ de elegidos con $p\ne q$, rechaza si $q\notin A_p$. Tiempo $\Oh(n^3)$.
\end{problema}

En \prob{Cocinero} el problema nuevo tiene un parámetro más (la disconformidad tolerada $d$); se fija $d=0$ y queda exactamente \npIS. Es el truco de <<restringir>>: si un caso particular de $X$ ya es $\NP$-completo, $X$ también lo es.\visto{Examen jul. 2022 Ej. 3: Cocinero, desde IS}

\begin{problema}{Cocinero}
\npdato{Problema} matriz simétrica $M$ de $n\times n$ con $M[i,j]=1$ si los ingredientes $i\ne j$ son incompatibles. La disconformidad de un plato $P$ es $\sum_{i,j\in P,\,i<j}M[i,j]$. Dados $k$ y $d$, ¿hay un plato con al menos $k$ ingredientes y disconformidad a lo sumo $d$?

\npdato{En $\NP$} certificado: el conjunto de ingredientes $t$, $|t|=\Oh(|s|)$. $B$ chequea $|t|\ge k$ y calcula la suma sobre los pares de $t$ en $\Oh(|t|^2)$; acepta si es $\le d$.

\npdato{Reducción $\npIS\redp\prob{Cocinero}$} $(G,k)\mapsto(M,k,0)$ con $M$ la matriz de adyacencia de $G$, en $\Oh(n^2)$. Disconformidad 0 significa que no hay pares incompatibles en el plato, es decir, un conjunto independiente.
\end{problema}

En \prob{RF} la dificultad está en traducir: cada arista pasa a ser un elemento de $X$ relacionado con sus dos extremos, y <<función parcial>> pasa a significar <<cada arista tiene a lo sumo un extremo elegido>>.\visto{Examen dic. 2022 Ej. 3: RF, desde IS}

\begin{problema}{Relación a función (RF)}
\npdato{Problema} una relación $R\subseteq X\times Y$ dada como matriz de ceros y unos (filas $Y$, columnas $X$) y un entero $k$: ¿hay un subconjunto $U\subseteq Y$ de al menos $k$ elementos tal que la restricción de $R$ a $U$ es una función parcial, o sea, cada $x\in X$ se relaciona con a lo sumo un elemento de $U$?

\npdato{En $\NP$} certificado: tira de $|Y|$ bits. $B$ chequea que tenga al menos $k$ unos y, columna por columna, que haya a lo sumo un 1 entre las filas elegidas. Tiempo $\Oh(|X|\,|Y|)$.

\npdato{Reducción $\npIS\redp\prob{rf}$} $Y=V$, $X=E$, y por cada arista $e=(u,w)$ se ponen en $R$ los pares $(e,u)$ y $(e,w)$; $k'=k$. Polinomial: la matriz tiene $n\cdot m$ celdas.

\npdato{Ida} si $U$ es independiente, cada arista tiene a lo sumo un extremo en $U$, así que cada $x=e$ se relaciona con a lo sumo un elemento de $U$.

\npdato{Vuelta (contrarrecíproco)} si $(G,k)$ es \npno, todo $U$ de $k$ vértices contiene una arista $(u,w)$, y entonces $x=(u,w)$ se relaciona con dos elementos de $U$: ningún $U$ sirve y $(R,k)$ es \npno.
\end{problema}

\subsection{Desde Clique}

Cuando lo que se busca es un subgrafo completo (o algo que contiene uno), conviene arrancar directamente de \npCL. La letra lo daba como sugerencia en los dos casos.

La reducción de \prob{Subgrafo Isomorfo} es la más corta del relevamiento: el grafo patrón es $K_k$. Lo interesante es el certificador, que recibe la biyección.\visto{Examen jul. 2024 Ej. 3: subgrafo isomorfo}

\begin{problema}{Subgrafo isomorfo}
\npdato{Problema} dados dos grafos $A$ y $G$ por matrices de adyacencia, ¿hay un subgrafo $B$ de $G$ isomorfo a $A$? ($A$ y $B$ son isomorfos si hay una biyección $f$ que preserva la adyacencia en ambos sentidos.)

\npdato{En $\NP$} certificado: lista de pares $(a,v)\in V_A\times V_G$ que define $f$. $B$ chequea que la lista tenga $|V_A|$ pares, que no se repitan vértices de $A$ ni de $G$ (así $f$ es función e inyectiva), y para cada par $(a,a')$ que $M_A[a,a']=M_G[f(a),f(a')]$. Tiempo $\Oh(|V_A|^2)$ más las búsquedas en la lista.

\npdato{Reducción $\npCL\redp\prob{Subgrafo Isomorfo}$} dada $(G,k)$ con $k\le|V|$, se construye $K$, el grafo completo de $k$ vértices, y se pregunta por $(K,G)$. Es $\Oh(k^2)=\Oh(|V|^2)$. Un clique de tamaño $k$ es un subgrafo isomorfo a $K$ y viceversa.
\end{problema}

\prob{Cometa} es el ejercicio con más trampa del relevamiento: la construcción natural (un solo camino colgado) no funciona, y la solución oficial lo dice explícitamente. Además hay que tratar aparte el caso $k>n$.\visto{Examen dic. 2023 Ej. 3: Cometa, desde Clique}

\begin{problema}{Cometa}
\npdato{Problema} un grafo de $2d$ vértices ($d\ge3$) es una cometa si $d$ vértices forman un clique y los otros $d$ forman un camino simple con un extremo conectado a un único vértice del clique. Dados $G$ (matriz de adyacencia) y $d\ge 3$, ¿hay un subgrafo de $G$ que es una cometa con al menos $2d$ vértices? Se sugiere \prob{Clique3}: \npCL\ con $k\ge 3$.

\npdato{En $\NP$} certificado: lista de $2d$ vértices, con el clique primero (el vértice de unión en la posición $d$) y el camino después. $B$ chequea que sean distintos, que los primeros $d$ sean adyacentes dos a dos y que cada par consecutivo desde la posición $d$ hasta la $2d$ sea adyacente. Son $\Oh(d^2)$ consultas a la matriz.

\npdato{Reducción $\prob{Clique3}\redp\prob{Cometa}$} si $k>n$, se responde \npno\ (o se devuelve una instancia \npno\ fija). Si no, $d=k$ y $G'$ es $G$ más, \emph{por cada vértice $v$}, un camino nuevo de $d$ vértices con un extremo unido a $v$ (caminos disjuntos). $G'$ tiene $n(k+1)$ vértices: polinomial porque $k\le n$.

\npdato{Ida} el clique de tamaño $k$ junto con el camino colgado de cualquiera de sus vértices es una cometa de $2d$ vértices.

\npdato{Vuelta} un vértice de un camino agregado tiene grado $\le 2$ y sus dos vecinos no son adyacentes, así que no está en ningún clique de tamaño $\ge3$. Como $d\ge3$, el clique de la cometa está entero dentro de $G$: $G$ tiene un clique de tamaño $k$.

\npdato{Lo que no sirve} colgar un único camino con un extremo adyacente a \emph{todos} los vértices de $G$: ese extremo junto con un clique de tamaño $d-1$ de $G$ forma un clique de tamaño $d$, y una instancia \npno\ de \npCL\ podría dar una cometa.
\end{problema}

\subsection{Desde Vertex Cover y Set Cover}

La señal es <<a lo sumo $k$ cosas que cubren todo>>. Hay dos formas del mismo problema y conviene distinguirlas para armar la transformación sin dudar:

\begin{itemize}
  \item \emph{forma \npSC}: se eligen \emph{grupos} y se quiere cubrir todos los \emph{elementos} (\prob{Test}: se eligen tests, se cubren líneas; \prob{Acceso}: se eligen lugares, se cubren dispositivos; \prob{uv}: se eligen caminos, se cubren vértices);
  \item \emph{forma \prob{Hitting Set}}: se eligen \emph{elementos} y se quiere tocar todos los \emph{grupos} (\prob{pcg}; \prob{Contratar}: se eligen aplicantes, cada tarea tiene su grupo de calificados).
\end{itemize}

Desde \npVC\ las dos son inmediatas: en la primera, cada vértice es un grupo (sus aristas incidentes) y cada arista es un elemento; en la segunda, cada vértice es un elemento y cada arista es un grupo de dos.

\prob{pcg} es literalmente \prob{Hitting Set}. La parte (c) solo pedía encadenar: $\npVC$ es $\NP$-completo, $\npVC\redp\prob{pcg}$ y se asume $\prob{pcg}\in\NP$.\visto{Examen feb. 2020 Ej. 1: PCG, desde VC}

\begin{problema}{Conjunto golpeador (PCG)}
\npdato{Problema} conjunto $A$, subconjuntos $B_1,\dots,B_m$ y $k$: ¿hay $S\subseteq A$ con $|S|\le k$ que contiene al menos un elemento de cada $B_i$?

\npdato{Reducción $\npVC\redp\prob{pcg}$} $A=V$; por cada arista $\{u,v\}$ un conjunto $B_i=\{u,v\}$; $k'=k$. Con listas de adyacencia, una pasada: para cada $v$ y cada vecino $w>v$ se agrega $\{v,w\}$ (la condición $w>v$ evita duplicados). Tiempo $\Oh(n+m)$.

\npdato{Ida y vuelta} $S$ golpea a todos los $B_i$ $\iff$ toda arista tiene un extremo en $S$ $\iff$ $S$ es un cubrimiento. El tamaño se conserva.
\end{problema}

En \prob{Test} los tests son vértices y las líneas son aristas; el test $v_i$ ejecuta las líneas de las aristas incidentes a $v_i$. Es la reducción del libro $\npVC\redp\npSC$ con otro vocabulario.\visto{Examen jul. 2021 Ej. 3: Test, desde VC}

\begin{problema}{Test}
\npdato{Problema} un programa de $N$ líneas y $M$ tests; el test $i$ ejecuta las líneas $L_i\subseteq\{1,\dots,N\}$. Dado $k$, ¿hay a lo sumo $k$ tests que ejecutan cada línea al menos una vez?

\npdato{Reducción $\npVC\redp\prob{Test}$} $N=m$ (una línea por arista), $M=n$ (un test por vértice), $L_i=\{j : e_j \text{ incide en } v_i\}$, $k'=k$. Tiempo $\Oh(nm)$.

\npdato{Ida} si $S$ cubre, los tests de $S$ ejecutan todas las líneas: la línea $j$ es la arista $e_j$, que tiene un extremo $v_{i_q}\in S$, y $j\in L_{i_q}$. \npdato{Vuelta} si $T$ ejecuta todas las líneas, para cada arista $e_j$ hay un test $i_q\in T$ con $j\in L_{i_q}$, o sea $e_j$ incide en $v_{i_q}$.

\npdato{En $\NP$} certificado: arreglo booleano de $M$ posiciones. $B$ chequea que haya a lo sumo $k$ elegidos y marca las líneas ejecutadas por los elegidos; acepta si quedan todas marcadas.
\end{problema}

\prob{Contratar} es la forma \prob{Hitting Set}: aplicantes $=$ vértices, tareas $=$ aristas, y los calificados para una tarea son los dos extremos de su arista. La letra sugería \npVC.\visto{2P 2022 Ej. 3: Contratar, desde VC}

\begin{problema}{Contratar}
\npdato{Problema} aplicantes $A=\{a_1,\dots,a_n\}$, tareas $t_1,\dots,t_m$ y, para cada tarea, el conjunto $C_i\subseteq A$ de calificados. Dado $k$, ¿se pueden contratar a lo sumo $k$ aplicantes de modo que cada tarea tenga al menos un calificado contratado?

\npdato{En $\NP$} certificado: arreglo $Z$ de $n$ bits. $B$ rechaza si $Z$ tiene más de $k$ unos; usa un arreglo $T$ de $m$ booleanos en \textit{false} y, para cada tarea $t_j$ y cada $a_i\in C_j$ con $Z[i]=1$, marca $T[j]$; acepta si todo $T$ quedó en \textit{true}. Tiempo $\Oh(nm)$.

\npdato{Reducción $\npVC\redp\prob{Contratar}$} un aplicante $a_i$ por vértice $v_i$, una tarea $t_i$ por arista $e_i=(v_j,v_k)$ con $C_i=\{a_j,a_k\}$, y $k'=k$. Tiempo $\Oh(n+m)$.

\npdato{Ida y vuelta} $S$ cubre la arista $e_j$ $\iff$ algún extremo de $e_j$ está en $S$ $\iff$ algún calificado para $t_j$ está contratado.
\end{problema}

\prob{Acceso} se reduce desde \npSC\ con la matriz de anchos de banda como matriz de pertenencia: $B(d,a_i)=1$ si $d\in S_i$ y 0 si no.\visto{2P 2023 Ej. 3: Acceso, desde Set Cover}

\begin{problema}{Acceso}
\npdato{Problema} dispositivos $D=\{d_1,\dots,d_n\}$, lugares $A=\{a_1,\dots,a_\ell\}$, una matriz $B$ de $n\times\ell$ con anchos de banda no negativos y un natural $k$. Un conjunto de lugares $S$ cubre a $d$ si $B(d,a)>0$ para algún $a\in S$. ¿Se puede cubrir todo $D$ con a lo sumo $k$ lugares?

\npdato{En $\NP$} certificado: vector de a lo sumo $k$ índices de lugares (más chico que la matriz). $B$ recorre las filas de la matriz y chequea que cada dispositivo tenga $B(d,a_i)>0$ para algún $a_i$ del certificado. Tiempo $\Oh(nk)$.

\npdato{Reducción $\npSC\redp\prob{Acceso}$} dada $(U,S_1,\dots,S_m,r)$: $D=U$, $\ell=m$, $k=r$, $B(d,a_i)=1$ si $d\in S_i$ y $0$ si no. Tiempo $\Oh(|U|\cdot m)$.

\npdato{Ida y vuelta} una subcolección de a lo sumo $r$ conjuntos cuya unión es $U$ corresponde exactamente a a lo sumo $k$ lugares que cubren todos los dispositivos.
\end{problema}

En \prob{uv} la letra guiaba paso a paso y pedía explícitamente identificar cuál problema es $X$ y cuál es $Y$ en $Y\redp X$: el origen $Y$ es \npSC, el destino $X$ es \prob{uv}. También daba como hipótesis la ida y la pertenencia a $\NP$.\visto{2P 2025 Ej. 2: UV, desde Set Cover}

\begin{problema}{Unión es V (UV)}
\npdato{Problema} un grafo $G=(V,E)$, un conjunto $\mathcal P=\{P_1,\dots,P_\ell\}$ de caminos simples de $G$ y un entero $k$: ¿hay $k$ caminos de $\mathcal P$ cuya unión de vértices es todo $V$?

\npdato{Reducción $\npSC\redp\prob{uv}$} dada $(U,\{S_1,\dots,S_\ell\},k)$: $V$ tiene un vértice $v_i$ por cada elemento $e_i\in U$; $G$ es el grafo completo sobre $V$ (así cualquier orden de vértices es un camino); por cada $S_i=\{e_{i_1},\dots,e_{i_h}\}$ se crea el camino $P_i=(v_{i_1},\dots,v_{i_h})$; el $k$ es el mismo. Si no se quiere el grafo completo, hay que agregar a mano las aristas de cada $P_i$. Es una copia: tiempo lineal en la instancia (más $\Oh(|U|^2)$ si se arma el completo).

\npdato{Vuelta (lo que se pedía)} si $k$ caminos $P_{i_1},\dots,P_{i_k}$ cubren $V$, para cada $e\in U$ su vértice $v$ está en algún $P_{i_j}$ y por construcción $e\in S_{i_j}$: los $S_{i_j}$ cubren $U$.

\npdato{Detalle} si \npSC\ pide <<a lo sumo $k$>> y \prob{uv} pide exactamente $k$, se supone $k\le\ell$ y, si hiciera falta, se agregan caminos de relleno; la solución oficial lo menciona en una nota al pie.
\end{problema}

\subsection{Desde SAT y 3-SAT}

La señal es que haya algo con dos estados por <<variable>> (fila, evento, ficha) y restricciones del tipo <<al menos uno de estos>>. La transformación pone una variable por cada objeto binario y una cláusula por cada restricción. Cuando el problema nuevo permite cláusulas de cualquier largo, se reduce desde \npSAT; cuando exige largo fijo, desde \npTSAT.

\prob{Solitaire} venía con la transformación dada: había que probar la vuelta, armar la reducción y concluir.\visto{2P 2019 Ej. 3: Solitaire, desde 3-SAT}

\begin{problema}{Solitaire}
\npdato{Problema} un tablero $n\times m$ con celdas vacías o con una ficha $X$ u $O$. Se gana si se puede quitar fichas de modo que en cada fila las que quedan sean todas del mismo tipo y cada columna tenga al menos una ficha.

\npdato{Transformación (dada)} desde una instancia de \npTSAT\ con variables $x_1..x_n$ y cláusulas $C_1..C_m$: tablero $n\times m$; si $x_i$ aparece en $C_j$ se pone $X$ en $(i,j)$, si aparece $\npneg{x_i}$ se pone $O$. Se recorren las cláusulas llenando una columna por cláusula: lineal.

\npdato{Vuelta} dada una configuración ganadora, $v(x_i)=1$ si la fila $i$ quedó con fichas $X$, $v(x_i)=0$ si quedó con $O$, arbitrario si quedó vacía (no hay ambigüedad porque cada fila es de un solo tipo). Cada columna $C_j$ conserva alguna ficha: si es $X$ en la fila $i$, el término $x_i$ es verdadero; si es $O$, el término $\npneg{x_i}$ es verdadero. Toda cláusula se satisface.

\npdato{Ida (se asumía)} con una asignación que satisface, en la fila $i$ se dejan solo las fichas del tipo que coincide con $v(x_i)$; cada columna conserva la ficha de su término verdadero.

\npdato{En $\NP$} certificado: lista de posiciones que se conservan. $B$ chequea que cada posición tenga ficha, que cada columna aparezca y que no haya dos fichas de distinto tipo en la misma fila. Con $R$ fichas, $\Oh(R^2)$ más los accesos al tablero.
\end{problema}

En \prob{casi-sat} se pide satisfacer todas menos una. El truco es agregar dos cláusulas contradictorias: exactamente una de ellas falla siempre, y por lo tanto las originales tienen que satisfacerse todas.\visto{2P 2020 Ej. 3: casiSAT, desde SAT}

\begin{problema}{casiSAT}
\npdato{Problema} variables $X$ y $k$ cláusulas: ¿hay una asignación que satisface al menos $k-1$ cláusulas?

\npdato{En $\NP$} certificado: la asignación. $B$ cuenta las cláusulas satisfechas recorriendo cada término una vez; acepta si son al menos $k-1$.

\npdato{Reducción $\npSAT\redp\prob{casi-sat}$} dada una instancia con variables $x_1..x_m$ y $r$ cláusulas, se agrega una variable nueva $z$ y dos cláusulas de un término: $(z)$ y $(\npneg z)$. Quedan $k=r+2$ cláusulas.

\npdato{Ida} la asignación que satisface las $r$ originales, con $z=1$, satisface $r+1=k-1$ cláusulas. \npdato{Vuelta} cualquier asignación falla en $(z)$ o en $(\npneg z)$; si satisface $k-1$, falla exactamente en una de esas dos y satisface las $r$ originales.
\end{problema}

\prob{Falsable} es el dual de \npSAT: una disyunción de conjunciones (forma normal disyuntiva) es falsa cuando todas las conjunciones son falsas, o sea, cuando cada una tiene un término falso.\visto{Examen feb. 2022 Ej. 2: Falsable, desde SAT}

\begin{problema}{Falsable}
\npdato{Problema} conjunciones $C_1,\dots,C_m$ (cada una $t_1\wedge\dots\wedge t_k$ con términos $x_i$ o $\npneg{x_i}$): ¿hay una asignación que hace falsa a $F=\bigvee_{i}C_i$?

\npdato{En $\NP$} certificado: la asignación. $B$ chequea que cada conjunción tenga algún término falso. Lineal.

\npdato{Reducción $\npSAT\redp\prob{Falsable}$} se niega la fórmula de \npSAT\ y se aplica De Morgan: cada cláusula $(t_1\vee\dots\vee t_r)$ pasa a la conjunción $(\npneg{t_1}\wedge\dots\wedge\npneg{t_r})$. Es cambiar el signo de cada término: lineal.

\npdato{Equivalencia} con la misma asignación, una cláusula original es verdadera si y solo si su conjunción negada es falsa. Todas las cláusulas son verdaderas si y solo si todas las conjunciones son falsas, es decir, $F$ es falsa.
\end{problema}

\prob{4-sat} parece trivial y no lo es: no se puede repetir una variable en una cláusula, así que <<agregar un cuarto término cualquiera>> no funciona si no se controla su valor. La solución oficial da dos reducciones.\visto{2P 2024 Ej. 2: 4-SAT, desde 3-SAT}

\begin{problema}{4-SAT}
\npdato{Problema} cláusulas de exactamente 4 términos sobre 4 variables distintas, representadas con matrices $V$ (índices) y $N$ (negaciones) de $k\times4$: ¿hay asignación que las satisface todas?

\npdato{En $\NP$} certificado: $n$ booleanos. $B$ evalúa las $k$ cláusulas en $\Oh(k)$. Como cada variable aparece en alguna cláusula, $n\le 4k$ y el certificado no es más largo que $N$.

\npdato{Reducción 1} se agregan variables nuevas $z,u,v,w$. Cada cláusula $C_i$ de \npTSAT\ pasa a $C_i\vee z$. Se agregan 8 cláusulas $(\ell_u\vee\ell_v\vee\ell_w\vee\npneg z)$, una por cada combinación de signos de $\ell_u\in\{u,\npneg u\}$, $\ell_v$, $\ell_w$. Para cualquier valor de $u,v,w$, una de esas 8 tiene sus tres primeros términos falsos, así que fuerzan $z=0$. Ida: misma asignación con $z=0$ y $u,v,w$ arbitrarios. Vuelta: $z=0$ forzado, luego cada $C_i\vee z$ se satisface por un término original.

\npdato{Reducción 2} cada cláusula $C_i$ se duplica: $C_i\vee x_{n+1}$ y $C_i\vee\npneg{x_{n+1}}$, con $x_{n+1}$ nueva ($2k$ cláusulas). Ida: misma asignación y $x_{n+1}$ arbitrario. Vuelta: cualquiera sea $x_{n+1}$, en una de las dos copias el cuarto término es falso y esa copia se satisface por un término de $C_i$.

\npdato{Tiempo} en las dos, agregar una columna a $V$ y $N$ y unas filas: lineal.
\end{problema}

En \prob{Al menos $k$} hay $k+1$ apostadores y se pide que acierten al menos $k$: otra vez <<todos menos uno>>, y otra vez el truco es agregar dos apuestas opuestas de las que exactamente una acierta.\visto{Examen jul. 2025 Ej. 3: Al menos k, desde SAT}

\begin{problema}{Al menos $k$}
\npdato{Problema} $n$ eventos (ganar $G_i$ o perder $P_i$) y $k+1$ apostadores, cada uno con apuestas sobre algunos eventos. ¿Hay resultados con los que al menos $k$ apostadores aciertan alguna apuesta? Se asume que hay un algoritmo polinomial que, dados la instancia y unos resultados, verifica si al menos $k$ aciertan.

\npdato{En $\NP$} el algoritmo dado es el certificador; el certificado es el conjunto de resultados, polinomial porque cada resultado corresponde a un evento con alguna apuesta. Hay que decir explícitamente las dos direcciones.

\npdato{Reducción $\npSAT\redp\prob{Al menos }k$} dada una instancia con $n$ variables y $c$ cláusulas: un evento por variable; un apostador por cláusula, con $G_i$ por cada término $x_i$ y $P_j$ por cada $\npneg{x_j}$; dos apostadores extra con $(G_n)$ y $(P_n)$. Hay $c+2$ apostadores y el umbral es $c+1$.

\npdato{Equivalencia} exactamente uno de los dos extra acierta. Por lo tanto aciertan al menos $c+1$ si y solo si aciertan los $c$ de las cláusulas. Con la correspondencia $x_i=1\iff$ se gana el evento $i$, el apostador de $C_h$ acierta si y solo si $C_h$ se satisface.

\npdato{Lo que no sirve} un apostador por cláusula sin los dos extra (umbral $c-1$): una instancia \npno\ de \npSAT\ que satisface $c-1$ cláusulas daría una instancia \npsi.
\end{problema}

\prob{nti4sat} (<<no todos iguales>>) pide en cada cláusula un término verdadero y uno falso. La reducción es agregar la misma variable nueva $s$ a todas las cláusulas; la clave de la vuelta es la simetría por complemento, que la letra daba como nota.\visto{Examen jul. 2026 Ej. 3: NTI4Sat, desde 3-SAT}

\begin{problema}{No todos iguales 4-SAT (NTI4Sat)}
\npdato{Problema} cláusulas de exactamente 4 términos sobre variables distintas. Una asignación satisface una cláusula si le da valor verdadero a algún término y falso a algún término. ¿Hay una asignación que satisface todas? Se asume que está en $\NP$.

\npdato{Reducción $\npTSAT\redp\prob{nti4sat}$} con una variable nueva $s$, cada cláusula $t_1\vee t_2\vee t_3$ pasa a $t_1\vee t_2\vee t_3\vee s$. Lineal.

\npdato{Ida} la asignación que satisface la instancia de \npTSAT, con $s=0$: cada cláusula tiene un término verdadero (de los originales) y uno falso ($s$).

\npdato{Vuelta} si una asignación satisface todas en el sentido <<no todos iguales>>, su complemento también (los verdaderos y los falsos se intercambian). Entre las dos, una tiene $s=0$; en esa, cada cláusula necesita un término verdadero, que tiene que ser uno de los tres originales. Restringida a las variables originales, satisface la instancia de \npTSAT.

\npdato{Lo que no sirve} agregar como cuarto término el opuesto de uno de los tres (repite variable: cláusula inválida, y además toda instancia sería \npsi), o agregar una fórmula en lugar de un término.
\end{problema}

\subsection{Desde camino o ciclo hamiltoniano}

La señal es <<recorrer todo>> o <<un camino o ciclo de largo al menos $k$>>. La reducción típica fija el parámetro para que <<largo $\ge k$>> obligue a pasar por todos los vértices.

El único ejercicio de este tipo daba camino hamiltoniano como $\NP$-completo y pedía la reducción y la pertenencia a $\NP$ por separado.\visto{2P 2021 Ej. 2: camino simple, desde CH}

\begin{problema}{Camino simple de largo al menos $k$ (CS)}
\npdato{Problema} dado un grafo $G$ y un natural $k$, ¿hay un camino simple (sin vértices repetidos) con al menos $k$ aristas?

\npdato{Reducción $\prob{ch}\redp\prob{cs}$} $G\mapsto(G,\,k=n-1)$. Copiar $G$ es lineal.

\npdato{Ida} un camino hamiltoniano no repite vértices y visita los $n$: es simple de largo $n-1$. \npdato{Vuelta} un camino simple de largo $n-1$ visita $n$ vértices distintos, o sea todos: es hamiltoniano.

\npdato{En $\NP$} certificado: lista de vértices $\pi$. $B$ rechaza si $\pi$ repite un vértice, si tiene $k$ o menos vértices, o si dos consecutivos no son adyacentes. Tiempo $\Oh(|\pi|^2+|\pi|m)$, y $|\pi|\le n$ porque no repite.
\end{problema}

Variantes que se resuelven igual: <<¿hay un ciclo de largo al menos $k$?>> (desde ciclo hamiltoniano con $k=n$), <<¿hay un camino hamiltoniano entre $s$ y $t$ dados?>>, o un \npTSP\ con restricciones extra (desde ciclo hamiltoniano con distancias 1 y 2).

\subsection{Desde Subset Sum, 3-coloreo y 3D-Matching}

En los parciales y exámenes relevados no se tomó ninguno con estos orígenes, pero están en el programa y son la elección natural para problemas numéricos o de etiquetado. Tres ejemplos cortos para practicar:

\begin{problema}{Partición (desde Subset Sum)}
\npdato{Problema} dados naturales $w_1,\dots,w_n$, ¿se pueden repartir en dos grupos de igual suma?

\npdato{Reducción} sea $S=\sum w_i$. Dada $(w,W)$ de \npSS, se agrega el número $\delta=|S-2W|$. Si $W\le S/2$, el total pasa a $2S-2W$ y cada mitad debe sumar $S-W$: un subconjunto que suma $W$ junto con $\delta$ suma $S-W$, y recíprocamente el lado que contiene a $\delta$, sin $\delta$, suma $W$. Si $W>S/2$, el total es $2W$ y el lado que no contiene a $\delta$ suma $W$. El número agregado tiene $\Oh(\log S)$ bits.
\end{problema}

\begin{problema}{4-coloreo (desde 3-coloreo)}
Se agrega un vértice nuevo adyacente a todos. En un 4-coloreo ese vértice tiene un color que no usa nadie más, así que el resto usa 3 colores; recíprocamente, un 3-coloreo se extiende dándole el cuarto color al vértice nuevo. En $\NP$: certificado, el color de cada vértice; se chequea cada arista.
\end{problema}

\begin{problema}{Mochila con valor exacto (desde Subset Sum)}
<<¿Hay un subconjunto de peso $\le W$ y valor exactamente $V$?>> Desde \npSS\ con $v_i=w_i$ y $V=W$. La lección general: si el problema nuevo tiene números y pide <<exactamente>>, \npSS\ casi siempre es el origen, y los números se copian (no se escriben en unario).
\end{problema}

% =====================================================================
\section{Mapa de lo que se preguntó}
% =====================================================================

Relevamiento de los segundos parciales (2008--2025) y exámenes (2009--2026) disponibles. Hasta 2016 los parciales y exámenes no incluían este tema; desde 2018 aparece en todos los segundos parciales y en buena parte de los exámenes, en general como el último ejercicio y con 17 a 40 puntos.

\begin{center}\small
\begin{tabularx}{\textwidth}{@{}>{\raggedright\arraybackslash}p{2.1cm} >{\raggedright\arraybackslash}X >{\raggedright\arraybackslash}p{3.3cm}@{}}
\toprule
\multicolumn{3}{@{}l}{\sans\bfseries Segundos parciales}\\
\midrule
Fuente & Qué se pidió & Reducido desde \\
\midrule
2P 2017 Ej. 4 & sí, no o no se sabe: Interval Scheduling y VC, IS e Interval Scheduling & conceptual \\
2P 2018 Ej. 3 & UOPI (independiente con pesos, valor $\ge t$) en $\NP$ y $\NP$-completo & \npIS \\
2P 2019 Ej. 3 & Solitaire: vuelta de una transformación dada, reducción y $\NP$-completo & \npTSAT \\
2P 2020 Ej. 3 & casiSAT (satisfacer al menos $k-1$ cláusulas) & \npSAT \\
2P 2021 Ej. 2 & camino simple de largo $\ge k$ & camino hamiltoniano \\
2P 2022 Ej. 3 & Contratar (aplicantes que cubren tareas) & \npVC \\
2P 2023 Ej. 3 & Acceso (puntos de acceso que cubren dispositivos) & \npSC \\
2P 2024 Ej. 2 & 4-SAT, con pauta de corrección detallada & \npTSAT \\
2P 2025 Ej. 1 & (e) transformación IS a VC; (f) definición de certificador eficiente & conceptual \\
2P 2025 Ej. 2 & UV (caminos cuya unión es $V$), guiado paso a paso & \npSC \\
\bottomrule
\end{tabularx}
\end{center}

\begin{center}\small
\begin{tabularx}{\textwidth}{@{}>{\raggedright\arraybackslash}p{2.1cm} >{\raggedright\arraybackslash}X >{\raggedright\arraybackslash}p{3.3cm}@{}}
\toprule
\multicolumn{3}{@{}l}{\sans\bfseries Exámenes}\\
\midrule
Fuente & Qué se pidió & Reducido desde \\
\midrule
Jul. 2018 P. 2 & $X\in\Pclass$, $Y\in\NP$, $Z$ $\NP$-completo: cuatro reducciones, ¿siempre, nunca o depende? & conceptual \\
Feb. 2019 Ej. 3 & definiciones de certificador eficiente, $\NP$ y $\NP$-completo & conceptual \\
Feb. 2020 Ej. 1 & PCG (conjunto golpeador) y definición de $\NP$ & \npVC \\
Feb. 2021 Ej. 3 & Red-Social (personas todas amigas entre sí) & \npIS\ (complemento) \\
Jul. 2021 Ej. 3 & Test (tests que ejecutan todas las líneas) & \npVC \\
Feb. 2022 Ej. 2 & Falsable (fórmula en forma disyuntiva que se puede hacer falsa) & \npSAT \\
Jul. 2022 Ej. 3 & Cocinero (ingredientes con disconformidad $\le d$) & \npIS\ ($d=0$) \\
Dic. 2022 Ej. 3 & RF (relación cuya restricción es función parcial) & \npIS \\
Dic. 2023 Ej. 3 & Cometa (clique de $d$ más camino de $d$) & \prob{Clique3} \\
Jul. 2024 Ej. 3 & subgrafo isomorfo & \npCL \\
Jul. 2025 Ej. 3 & Al menos $k$ (apostadores que aciertan) & \npSAT \\
Dic. 2025 Ej. 1 & (e) certificado de SAT; (f) definición de $\NP$-completo y cómo usarla & conceptual \\
Jul. 2026 Ej. 1 & (f) certificado de SAT & conceptual \\
Jul. 2026 Ej. 3 & NTI4Sat (no todos iguales, 4 términos) & \npTSAT \\
\bottomrule
\end{tabularx}
\end{center}

Lo que se repite: el origen es \npIS, \npVC/\npSC\ o \npSAT/\npTSAT\ en 15 de los 18 ejercicios de reducción (los otros tres salen de \npCL\ o de camino hamiltoniano); la transformación es casi siempre una copia con otro nombre; y lo que más puntaje cuesta es la redacción completa del certificador y la vuelta de la equivalencia.
